\documentclass{article}
\usepackage{amsmath}
\usepackage{algorithm}
\usepackage{algorithmic}
\usepackage[utf8]{inputenc}
\usepackage[spanish]{babel}

\title{Ejemplos de uso de comPyLaTex}
\author{}
\date{}

\begin{document}

\maketitle


\section{C\'alculo de ra\'{\i}ces}

\begin{equation}
f(x) = x \cdot \ln(2)
\label{eq:calc:f}
\end{equation}

\begin{equation}
a_1 = -1
\label{eq:calc:a1}
\end{equation}

\begin{equation}
a_2 = 10
\label{eq:calc:a2}
\end{equation}

\begin{equation}
    tol_{dic} = 1e-6
    \label{eq:calc:toldic}
\end{equation}

\begin{algorithm}
\caption{M\'etodo de dicotom\'{\i}a}
\begin{algorithmic}
    \label{alg:calc:a3}
    \WHILE{$a_2 - a_1 > tol_{dic}$}
        \STATE $a_3 = (a_1 + a_2)/2$
        \IF{$f(a_1) \cdot f(a_3) < 0$}
            \STATE $a_2 = a_3$
        \ELSE
            \STATE $a_1 = a_3$
        \ENDIF
    \ENDWHILE
    \RETURN $a_3$
\end{algorithmic}
\end{algorithm}


\begin{equation}
    a_3 = 
    \label{eq:res:a3}
\end{equation}


\section{C\'alculo de soluciones de $x^a + y^b = z^c$}

Valores $a$, $b$, y $c$
\begin{equation}
a = 1
\label{eq:calc:a}
\end{equation}

\begin{equation}
b = 2
\label{eq:calc:b}
\end{equation}

\begin{equation}
c = 2
\label{eq:calc:c}
\end{equation}

Valor m\'aximo para las soluciones ($x, y, z \in \{1,2,..max\}$):
\begin{equation}
max = 10
\label{eq:calc:max}
\end{equation}



\begin{algorithm}
\caption{N\'umero de soluciones}
\begin{algorithmic}
\label{alg:calc:sols}
    \STATE sols = 0
    \STATE z = 1
    \WHILE{$z \leq max$}
        \STATE y = 1
        \WHILE{$y \leq max$}
            \STATE x = 1
            \WHILE{$x \leq max$}
                \IF{$x^a + y^b = z^c$}
                    \STATE sols = sols + 1
                \ENDIF
                \STATE x = x + 1
            \ENDWHILE
            \STATE y = y + 1
        \ENDWHILE
        \STATE z = z + 1
    \ENDWHILE
    \RETURN sols
\end{algorithmic}
\end{algorithm}

\begin{equation}
    sols = 
    \label{eq:res:sols}
\end{equation}


\section{Ejemplo con ELSIFs}

\begin{algorithm}
\caption{Ejemplo de un algoritmo con ELSIFs anidados}
\begin{algorithmic}
\label{alg:calc:j}
        \STATE $i = 10$
        \STATE $j = \sqrt{i^2 + 5}$
        \IF{$i \ge 5$}
            \STATE $i = i - 1$
            \WHILE{$j \le 20$}
                \STATE $j = j + \log(i)$
            \ENDWHILE
        \ELSIF{$i \le 3$}
            \STATE $i = i + 2$
        \ENDIF
        \RETURN $j$
\end{algorithmic}
\end{algorithm}



\begin{equation}
    j = 
    \label{eq:res:j}
\end{equation}

\section{Prueba con ra\'{\i}ces c\'ubicas}
\begin{equation}
    a = \sqrt[3]{8}
    \label{eq:calc:a}
\end{equation}

\begin{equation}
    a = 
    \label{eq:res:a}
\end{equation}


% ============================================================
% SECCI\'ON 1: C\'ALCULO — Newton-Raphson
% ============================================================

\section{C\'alculo: M\'etodo de Newton-Raphson}

El m\'etodo de Newton-Raphson permite aproximar ra\'{\i}ces de funciones diferenciables
mediante la iteraci\'on $x_{n+1} = x_n - \frac{f(x_n)}{f'(x_n)}$.
Aplicamos el m\'etodo a $f(x) = x^3 - 2$, cuya ra\'{\i}z es $\sqrt[3]{2} \approx 1.259921$.

\begin{equation}
    f(x) = x^3 - 2
    \label{eq:calc:f}
\end{equation}

\begin{equation}
    df(x) = 3 \cdot x^2
    \label{eq:calc:df}
\end{equation}

\begin{equation}
    x_0 = 1.5
    \label{eq:calc:x0}
\end{equation}

\begin{equation}
    tol_{NR} = 1e-6
    \label{eq:calc:tolNR}
\end{equation}

\begin{algorithm}
\caption{M\'etodo de Newton-Raphson para $f(x) = x^3 - 2$}
\begin{algorithmic}
\label{alg:calc:x}
    \STATE $x = x_0$
    \WHILE{$\left|f(x)\right| > tol_{NR}$}
        \STATE $x = x - f(x) / df(x)$
    \ENDWHILE
    \RETURN $x$
\end{algorithmic}
\end{algorithm}

La ra\'{\i}z aproximada de $f(x) = x^3 - 2$ obtenida por Newton-Raphson es:

\begin{equation}
    x =
    \label{eq:res:x}
\end{equation}

% ============================================================
% SECCI\'ON 2: \'ALGEBRA — M\'etodo de Jacobi
% ============================================================

\section{\'Algebra: M\'etodo de Jacobi}

El m\'etodo de Jacobi es un m\'etodo iterativo para la resoluci\'on de sistemas de
ecuaciones lineales. Aplicamos el m\'etodo al sistema:
\[
    \left \{
        \begin{array}{rr}
        5x-y+z & = 10 \\
        2x+8y-z & = 11 \\
        -x+y+4z & = 3
        \end{array}
    \right .
\]

cuya soluci\'on exacta es $(x, y, z) = (2,1,1)$.

Las ecuaciones de actualizaci\'on de Jacobi para este sistema son:
\[
    x_{n+1} = \frac{10 + y_n - z_n}{5}, \quad y_{n+1} = \frac{11 - 2\cdot x_n + z_n}{8}, \quad z_{n+1} = \frac{3 + x_n - y_n}{4}
\]

\begin{equation}
    x_j = 0
    \label{eq:calc:xj}
\end{equation}

\begin{equation}
    y_j = 0
    \label{eq:calc:yj}
\end{equation}

\begin{equation}
    z_j = 0
    \label{eq:calc:zj}
\end{equation}

\begin{equation}
    tol_{J} = 1e-6
    \label{eq:calc:tolJ}
\end{equation}

\begin{algorithm}
\caption{M\'etodo de Jacobi para $5x-y+z = 10$, $2x+8y-z = 11$, $-x+y+4z = 3$}
\begin{algorithmic}
\label{alg:calc:x,y,z}
    \STATE $x = (10 + y_j - z_j)/5$
    \STATE $y = (11 - 2\cdot x_j + z_j)/8$
    \STATE $z = (3 + x_j - y_j)/4$
    \WHILE{$\left|x - x_j\right| > tol_{J}$}
        \STATE $x_j = x$
        \STATE $y_j = y$
        \STATE $z_j = z$
        \STATE $x = (10 + y_j - z_j)/5$
        \STATE $y = (11 - 2\cdot x_j + z_j)/8$
        \STATE $z = (3 + x_j - y_j)/4$
    \ENDWHILE
    \RETURN $x$, $y$, $z$
\end{algorithmic}
\end{algorithm}

La soluci\'on aproximada del sistema obtenida por Jacobi es:

\begin{equation}
    x =
    \label{eq:res:x}
\end{equation}

\begin{equation}
    y =
    \label{eq:res:y}
\end{equation}

\begin{equation}
    z =
    \label{eq:res:z}
\end{equation}

% ============================================================
% SECCI\'ON 3: AN\'ALISIS — M\'etodo de Euler para EDOs
% ============================================================

\section{An\'alisis: M\'etodo de Euler para ecuaciones diferenciales}

El m\'etodo de Euler expl\'{\i}cito aproxima la soluci\'on de un problema de valor inicial
$y' = f(t, y)$, $y(t_0) = y_0$, mediante la iteraci\'on:
\[
    y_{n+1} = y_n + h \cdot f(t_n, y_n)
\]
Aplicamos el m\'etodo a $y' = y$, $y(0) = 1$, cuya soluci\'on exacta es $y(t) = e^t$.
Evaluamos en $t = 1$, donde la soluci\'on exacta vale $e \approx 2.718282$.

\begin{equation}
    g(t, y) = y
    \label{eq:calc:g}
\end{equation}

\begin{equation}
    y_0 = 1
    \label{eq:calc:y0}
\end{equation}

\begin{equation}
    h = 0.01
    \label{eq:calc:h}
\end{equation}

\begin{equation}
    T = 1
    \label{eq:calc:T}
\end{equation}

\begin{algorithm}
\caption{M\'etodo de Euler para $y' = y$, $y(0) = 1$}
\begin{algorithmic}
\label{alg:calc:y}
    \STATE $t = 0$
    \STATE $y = y_0$
    \WHILE{$t < T$}
        \STATE $y = y + h \cdot g(t, y)$
        \STATE $t = t + h$
    \ENDWHILE
    \RETURN $y$
\end{algorithmic}
\end{algorithm}

La aproximaci\'on de $e = y(1)$ obtenida por el m\'etodo de Euler con paso $h = 0.01$ es:

\begin{equation}
    y =
    \label{eq:res:y}
\end{equation}

% ============================================================
% SECCI\'ON 4: ESTAD\'ISTICA — N\'umero combinatorio
% ============================================================

\section{Estad\'{\i}stica: N\'umero combinatorio}

El n\'umero combinatorio $\binom{n}{k}$ cuenta el n\'umero de subconjuntos de tama\~no
$k$ de un conjunto de $n$ elementos. Se calcula mediante:
\[
    \binom{n}{k} = \frac{n!}{k! \cdot (n-k)!}
\]
Calculamos $\binom{10}{3} = 120$ mediante el c\'alculo iterativo de factoriales.

\begin{equation}
    n = 10
    \label{eq:calc:n}
\end{equation}

\begin{equation}
    k = 3
    \label{eq:calc:k}
\end{equation}

\begin{algorithm}
\caption{C\'alculo de $\binom{n}{k}$ mediante factoriales iterativos}
\begin{algorithmic}
    \label{alg:calc:nc}
    \STATE $fn = 1$
    \STATE $i = 1$
    \WHILE{$i \leq n$}
        \STATE $fn = fn \cdot i$
        \STATE $i = i + 1$
    \ENDWHILE
    \STATE $fk = 1$
    \STATE $i = 1$
    \WHILE{$i \leq k$}
        \STATE $fk = fk \cdot i$
        \STATE $i = i + 1$
    \ENDWHILE
    \STATE $fnk = 1$
    \STATE $i = 1$
    \WHILE{$i \leq n - k$}
        \STATE $fnk = fnk \cdot i$
        \STATE $i = i + 1$
    \ENDWHILE
    \STATE $nc = fn / (fk \cdot fnk)$
    \RETURN $nc$
\end{algorithmic}
\end{algorithm}

El n\'umero combinatorio $\binom{10}{3}$ calculado por el programa es:

\begin{equation}
    \binom{10}{3} =
    \label{eq:res:nc}
\end{equation}

% ============================================================
% SECCI\'ON 5: BIOLOG\'IA MATEM\'ATICA — Lotka-Volterra
% ============================================================

\section{Biolog\'{\i}a matem\'atica: Sistema de Lotka-Volterra}

El sistema de Lotka-Volterra modela la din\'amica de poblaciones depredador-presa:
\[
\left \{
\begin{array}{rrrr}
    \frac{dx}{dt} & = &\alpha x - \beta x y \\
    \quad \frac{dy}{dt} & = &\delta x y - \gamma y
\end{array}
\right .
\]
donde $x$ es la poblaci\'on de presas, $y$ la de depredadores, y $\alpha$, $\beta$,
$\delta$, $\gamma$ son par\'ametros del modelo. Aplicamos Euler expl\'{\i}cito con los
par\'ametros cl\'asicos $\alpha = 1.0$, $\beta = 0.1$, $\delta = 0.075$, $\gamma = 1.5$,
condiciones iniciales $x(0) = 10$, $y(0) = 5$, paso $h = 0.01$ y tiempo final $T = 10$.

\begin{equation}
    \alpha = 1.0
    \label{eq:calc:alpha}
\end{equation}

\begin{equation}
    \beta = 0.1
    \label{eq:calc:beta}
\end{equation}

\begin{equation}
    \delta = 0.075
    \label{eq:calc:delta}
\end{equation}

\begin{equation}
    \gamma = 1.5
    \label{eq:calc:gamma}
\end{equation}

\begin{equation}
    x_{LV} = 10.0
    \label{eq:calc:xlv}
\end{equation}

\begin{equation}
    y_{LV} = 5.0
    \label{eq:calc:ylv}
\end{equation}

\begin{equation}
    h_{LV} = 0.01
    \label{eq:calc:hlv}
\end{equation}

\begin{equation}
    T_{LV} = 10
    \label{eq:calc:Tlv}
\end{equation}

\begin{algorithm}
\caption{Euler expl\'{\i}cito para el sistema de Lotka-Volterra}
\begin{algorithmic}
\label{alg:calc:xlv}
    \STATE $t_{LV} = 0$
    \WHILE{$t_{LV} < T_{LV}$}
        \STATE $dxdt = \alpha \cdot x_{LV} - \beta \cdot x_{LV} \cdot y_{LV}$
        \STATE $dydt = \delta \cdot x_{LV} \cdot y_{LV} - \gamma \cdot y_{LV}$
        \STATE $x_{LV} = x_{LV} + h_{LV} \cdot dxdt$
        \STATE $y_{LV} = y_{LV} + h_{LV} \cdot dydt$
        \STATE $t_{LV} = t_{LV} + h_{LV}$
    \ENDWHILE
    \RETURN $x_{LV}$
\end{algorithmic}
\end{algorithm}

Las poblaciones aproximadas en $t = 10$ obtenidas por el m\'etodo de Euler son:

\begin{equation}
    x_{LV} =
    \label{eq:res:xlv}
\end{equation}

\begin{equation}
    y_{LV} =
    \label{eq:res:ylv}
\end{equation}


\end{document}